\documentclass{bmstu}

\usepackage{svg}
\usepackage{adjustbox}

\usetikzlibrary{positioning,shapes.geometric}

\addbibresource{references.bib}

\begin{document}

\makereporttitle
{Информатика и системы управления}
{Программное обеспечение ЭВМ и информационные технологии}
{лабораторной работе №2}
{Анализ алгоритмов}
{Сравнительный анализ ресурсной эффективности алгоритмов}
{} % Вариант
{Кузнецов М. А./ИУ7-52Б} % Номер группы/ФИО студента (если несколько - разделить запятой)
{Строганов Д. В.}

\maketableofcontents

\chapter*{ВВЕДЕНИЕ}
\addcontentsline{toc}{chapter}{ВВЕДЕНИЕ}

Цель работы --- выполнение сравнительного анализа ресурсной
эффективности алгоритмов умножения матриц по трудоемкости, памяти и времени
выполнения их реализаций.

Задачи работы:

\renewcommand{\labelitemi}{---}
\begin{itemize}
	\item описать решаемую задачу согласно варианту;
	\item разработать алгоритмы решения задачи согласно варианту ---
	      стандартный алгоритм умножения матриц, алгоритм Винограда и три
	      оптимизированые версии алгоритма Винограда;
	\item описать разработанные алгоритмы схемами алгоритмов;
	\item описать требования к программному обеспечению;
	\item описать средства реализации алгоритмов;
	\item реализовать разработанные алгоритмы;
	\item выполнить тестирование реализации алгоритмов;
	\item Описать модель вычислений для оценки трудоёмкости алгоритмов и
	      выполнить теоретический расчёт трудоёмкости;
	\item описать инструменты замера процессорного времени выполнения
	      реализаций алгоритмов;
	\item выполнить массированные замеры процессорного времени
	      выполнения реализаций алгоритмов в зависимости от линейного
	      размера квадратных матриц, отдельно для чётных и нечётных
	      размеров;
	\item выполнить сравнительный анализ алгоритмов по трудоёмкости и по
	      данным выполненных замеров;
\end{itemize}

\chapter{Аналитическая часть}

\section{Постановка задачи}
Даны матрица $A$ размера $n × m$ и матрица $B$ размера $m × q$ с
целочисленными элементами. Требуется вычислить их произведение —
матрицу $C = A · B$ размера $n × q$, элементы которой определяются
формулой:

\begin{equation}
	c_{i,j} = \sum_{k=0}^{m-1} a_{i,k} \cdot b_{k,j}, \quad i = \overline{0, n - 1}, \quad j = \overline{0, q - 1},
\end{equation}

где $c_{i,j}$ --- элемент матрицы $C$ в строке $i$ и столбце $j$;

$k$ --- индекс суммирования;

$m$ --- число столбцов матрицы $A$, равное числу строк матрицы $B$;

$a_{i,k}$, $b_{k,j}$ --- элементы матриц $A$ и $B$;

$n$ --- число строк матрицы $A$;

$q$ --- число столбцов матрицы $B$.

Произведение определено только в том случае, когда число столбцов матрицы $A$ равно числу строк матрицы $B$.

Для решения задачи рассматриваются следующие алгоритмы:
\begin{itemize}
	\item стандартный алгоритм умножения матриц;
	\item алгоритм Винограда;
	\item оптимизированная согласно варианту реализации алгоритма Винограда, в которых выполняется кэширование неизменяющихся величин, а в залоговке цикла применяется операция $i += 2$.
\end{itemize}

\section{Стандартный алгоритм умножения матриц}
Стандартный алгоритм вычисляет каждый элемент результирующей матрицы как
скалярное произведение строки матрицы $A$ и столбца матрицы $B$. Для всех
$i \in \{0,\ldots,n-1\}$ и $j \in \{0,\ldots,q-1\}$ сначала
инициализируется $c_{i,j}=0$, после чего выполняется накопление суммы:
\begin{equation}
	c_{i,j} = c_{i,j} + a_{i,k} \cdot b_{k,j}, \quad
	k=0,\ldots,m-1.
\end{equation}
Таким образом, алгоритм содержит три вложенных цикла: внешний цикл перебирает
строки матрицы $A$, следующий цикл --- столбцы матрицы $B$, а внутренний цикл
формирует одно значение результирующей матрицы.

\section{Алгоритм Винограда}
Алгоритм Винограда преобразует сумму произведений, объединяя соседние
слагаемые попарно. Для строки $i$ матрицы $A$ и столбца $j$ матрицы $B$
предварительно вычисляются множители:
\begin{equation}
	R_i = \sum_{k=0}^{\lfloor m/2\rfloor-1}
	a_{i,2k}\cdot a_{i,2k+1},
	\label{eq:winograd-precompute-1}
\end{equation}

\begin{equation}
	C_j = \sum_{k=0}^{\lfloor m/2\rfloor-1}
	b_{2k,j}\cdot b_{2k+1,j}.
	\label{eq:winograd-precompute-2}
\end{equation}
Для чётного $m$ элемент результирующей матрицы вычисляется по формуле
\begin{equation}
	c_{i,j} =
	-R_i-C_j+
	\sum_{k=0}^{\lfloor m/2\rfloor-1}
	\left(a_{i,2k}+b_{2k+1,j}\right)
	\cdot
	\left(a_{i,2k+1}+b_{2k,j}\right).
\end{equation}

Если $m$ нечётно, к правой части последней формулы дополнительно прибавляется
последнее непарное произведение:
\begin{equation}
	c_{i,j} = c_{i,j} + a_{i,m-1}\cdot b_{m-1,j}.
\end{equation}
Вычисление $R_i$ и $C_j$ до основного цикла позволяет не повторять
неизменяющиеся произведения для каждой пары $i,j$.

\section{Вывод}
Для матриц размеров $n\times m$ и $m\times q$ стандартный алгоритм выполняет
$nmq$ умножений и $nq(m-1)$ сложений. Его временная сложность составляет
$O(nmq)$, а дополнительная память без учёта матрицы результата имеет
сложность $O(1)$.

Алгоритм Винограда также имеет временную сложность $O(nmq)$, поскольку
число обрабатываемых пар индексов пропорционально $m/2$, а все пары
строк и столбцов перебираются. При этом предварительное вычисление $R_i$ и
$C_j$ уменьшает число умножений, добавляя операции сложения и вычитания.
Дополнительная память алгоритма составляет $O(n+q)$ для массивов
предварительно вычисленных значений; память под результирующую матрицу не
учитывается. При квадратных матрицах оба алгоритма имеют сложность
$O(n^3)$, но различаются количеством базовых операций и объёмом
вспомогательной памяти.


\chapter{Конструкторская часть}

\section{Описание алгоритмов}
Для каждого алгоритма приняты следующие обозначения:
\begin{itemize}
	\item $a$ — матрица $A$ размера $n × m$;
	\item $b$ — матрица $B$ размера $m × q$;
	\item $c$ — результирующая матрица размера $n × q$, изначально
	      заполненная нулями.
\end{itemize}

Для алгоритмов Винограда дополнительно определены:
\begin{itemize}
	\item $R$ — массив из $n$ элементов, вычисляемый по формуле~\ref{eq:winograd-precompute-1};
	\item $C$ — массив из $q$ элементов, вычисляемый по формуле~\ref{eq:winograd-precompute-2}.
\end{itemize}

\subsection{Стандартный алгоритм умножения матриц}
Схема стандартного алгоритма умножения матриц приведена на рисунке~\ref{fig:standart_matrix_multiplication}.
\begin{figure}[H]
	\centering
	\begin{adjustbox}{max width=\textwidth, max height=0.8\textheight}
		\includesvg[inkscapelatex=false]
		{../../diagrams/lab_2/standart_matrix_multiplication}
	\end{adjustbox}
	\caption{Схема стандартного алгоритма умножения матриц}
	\label{fig:standart_matrix_multiplication}
\end{figure}

\subsection{Алгоритм Винограда}
Схема алгоритма Винограда для умножения матриц приведена на рисунке~\ref{fig:vinograd_matrix_multiplication}.
\begin{figure}[H]
	\centering
	\begin{adjustbox}{max width=\textwidth, max height=0.8\textheight}
		\includesvg[inkscapelatex=false]
		{../../diagrams/lab_2/vinograd_matrix_multiplication}
	\end{adjustbox}
	\caption{Схема алгоритма Винограда умножения матриц}
	\label{fig:vinograd_matrix_multiplication}
\end{figure}

\subsection{Оптимизированный алгоритм Винограда}
Согласно варианту к алгоритму Винограда применяются две оптимизации:
\begin{itemize}
	\item Кеширование неизменяющихся, но многократно вычисляемых величин ---
	      $m // 2$ в условии циклов и $m \% 2$ в проверке на нечетность.
	\item Изменения шага цикла на $i += 2$ для отказа от умножения на $2$ при
	      вычислении позиции элемента в матрице.
\end{itemize}

Схема оптимизированного алгоритма Винограда для умножения матриц приведена на рисунке~\ref{fig:optimized_vinograd_matrix_multiplication}.
\begin{figure}[H]
	\centering
	\begin{adjustbox}{max width=\textwidth, max height=0.8\textheight}
		\includesvg[inkscapelatex=false]
		{../../diagrams/lab_2/optimized_vinograd_matrix_multiplication}
	\end{adjustbox}
	\caption{Схема алгоритма Винограда умножения матриц}
	\label{fig:optimized_vinograd_matrix_multiplication}
\end{figure}

\section{Требования к программному обеспечению}
К программе предъявляются следующие требования:
\begin{itemize}
	\item программа загружает матрицы A и B из текстовых файлов, пути
	      к которым вводит пользователь; одна строка файла соответствует
	      одной строке матрицы, элементы — целые числа, разделённые
	      пробелами;
	\item при загрузке программа проверяет, что все элементы являются
	      целыми числами, а все строки матрицы содержат одинаковое число
	      элементов; при ошибке выводится сообщение, а ранее загруженная
	      матрица сохраняется;
	\item перед умножением программа проверяет, что число столбцов
	      матрицы A равно числу строк матрицы B;
	\item программа циклически выводит меню со следующими пунктами:
	      \begin{itemize}
		      \item выполнение умножения выбранным алгоритмом;
		      \item выполнение умножения всеми алгоритмами;
		      \item вывод загруженных матриц;
		      \item выбор алгоритма;
		      \item загрузка матрицы A из файла;
		      \item загрузка матрицы B из файла;
		      \item замеры процессорного времени;
		      \item завершение работы.
	      \end{itemize}
	\item  в режиме замеров программа для каждого линейного размера N
	      генерирует случайные квадратные матрицы, выполняет их
	      умножение каждым алгоритмом, выводит среднее процессорное
	      время на экран и сохраняет результаты в файлы формата CSV, по
	      одному на алгоритм.
	\item программа должна быть реализована на языке программирования со статической типизацией.
\end{itemize}

\chapter{Технологическая часть}

\section{Средства реализации}
В качестве языка программирования со статической типизацией выбран C++
(стандарт C++20) \cite{cpp20}. % TODO кратко перечисли использованные возможности стандартной библиотеки

Для сборки исполняемых файлов предполагается использовать компилятор GCC
с флагами \texttt{-std=c++20 -Wall -Wextra -O0}.

\section{Реализация алгоритмов}
Объявления разработанных функций вынесены в заголовочные файлы
\texttt{matrix.hpp} и \texttt{benchmark\_support.hpp} (листинги~\ref{lst:matrix-hpp}
и \ref{lst:benchmark-support-hpp}), а определения --- в файлы \texttt{matrix.cpp},
\texttt{benchmark\_support.cpp} и \texttt{benchmark.cpp} (листинги~\ref{lst:matrix-cpp},
\ref{lst:benchmark-support-cpp} и \ref{lst:benchmark-cpp}). Программа
с интерактивным меню реализована в файле \texttt{main.cpp} (Листинг~\ref{lst:main-cpp}).

\section{Тестирование реализации алгоритмов}
Для проверки корректности реализации алгоритмов был реализован модуль
тестирования \texttt{test.cpp} (листинг~\ref{lst:test-cpp}). Каждый набор входных данных проверяется
с помощью стандартного алгоритма, алгоритма Винограда и оптимизированного
алгоритма Винограда. Результат работы алгоритмов сравнивается с заранее
известной матрицей.

\begin{table}[H]
	\centering
	\caption{Описание тестов}
	\label{tab:tests}
	\begin{tabularx}{\textwidth}{|X|X|X|X|}
		\hline
		Описание теста & Входные данные                               & Ожидаемый результат &
		Полученный результат                                                                  \\ \hline
		Умножение прямоугольных матриц
		               & $\begin{pmatrix}1&2&3\\4&5&6\end{pmatrix}$,
		$\begin{pmatrix}7&8\\9&10\\11&12\end{pmatrix}$
		               & $\begin{pmatrix}58&64\\139&154\end{pmatrix}$
		               & $\begin{pmatrix}58&64\\139&154\end{pmatrix}$                         \\ \hline
		Нечётный внутренний размер матриц
		               & $\begin{pmatrix}2&-1&4\\0&3&5\end{pmatrix}$,
		$\begin{pmatrix}1&2\\-2&0\\3&4\end{pmatrix}$
		               & $\begin{pmatrix}16&20\\9&20\end{pmatrix}$
		               & $\begin{pmatrix}16&20\\9&20\end{pmatrix}$                            \\ \hline
		Умножение на единичную матрицу слева
		               & $\begin{pmatrix}1&0\\0&1\end{pmatrix}$,
		$\begin{pmatrix}2&-1&4\\0&3&5\end{pmatrix}$
		               & $\begin{pmatrix}2&-1&4\\0&3&5\end{pmatrix}$
		               & $\begin{pmatrix}2&-1&4\\0&3&5\end{pmatrix}$                          \\ \hline
		Умножение на единичную матрицу справа
		               & $\begin{pmatrix}2&-1&4\\0&3&5\end{pmatrix}$,
		$\begin{pmatrix}1&0&0\\0&1&0\\0&0&1\end{pmatrix}$
		               & $\begin{pmatrix}2&-1&4\\0&3&5\end{pmatrix}$
		               & $\begin{pmatrix}2&-1&4\\0&3&5\end{pmatrix}$                          \\ \hline
		Умножение нулевой матрицы слева
		               & $\begin{pmatrix}0&0&0\\0&0&0\end{pmatrix}$,
		$\begin{pmatrix}7&8\\9&10\\11&12\end{pmatrix}$
		               & $\begin{pmatrix}0&0\\0&0\end{pmatrix}$
		               & $\begin{pmatrix}0&0\\0&0\end{pmatrix}$                               \\ \hline
		Умножение нулевой матрицы справа
		               & $\begin{pmatrix}1&2&3\\4&5&6\end{pmatrix}$,
		$\begin{pmatrix}0&0\\0&0\\0&0\end{pmatrix}$
		               & $\begin{pmatrix}0&0\\0&0\end{pmatrix}$
		               & $\begin{pmatrix}0&0\\0&0\end{pmatrix}$                               \\ \hline
		Умножение обратных матриц
		               & $\begin{pmatrix}1&1\\0&1\end{pmatrix}$,
		$\begin{pmatrix}1&-1\\0&1\end{pmatrix}$
		               & $\begin{pmatrix}1&0\\0&1\end{pmatrix}$
		               & $\begin{pmatrix}1&0\\0&1\end{pmatrix}$                               \\ \hline
		Умножение матриц размера $1\times1$
		               & $\begin{pmatrix}-1\end{pmatrix}$,
		$\begin{pmatrix}-1\end{pmatrix}$
		               & $\begin{pmatrix}1\end{pmatrix}$
		               & $\begin{pmatrix}1\end{pmatrix}$                                      \\ \hline
	\end{tabularx}
\end{table}

Все тесты были пройдены.

\chapter{Исследовательская часть}

\section{Оценка временной эффективности}
% Временная эффективность характеризуется числом сравнений искомого числа с
% элементами массива.
% \subsection{Линейный поиск}
% В лучшем случае искомый элемент находится на первой позиции массива, поэтому
% требуется одно сравнение: $T_{\min}(N)=1=O(1)$. В худшем случае элемент
% находится на последней позиции или отсутствует, поэтому выполняется $N$
% сравнений: $T_{\max}(N)=N=O(N)$.

% \subsection{Бинарный поиск, нерекурсивная реализация}
% Для варианта 1 в лучшем случае выполняется одно сравнение с центральным
% элементом: $T_{\min}(N)=1=O(1)$. В худшем случае выполняется
% $\lceil\log_2(N+1)\rceil$ шагов по два сравнения:
% $T_{\max}(N)=2\lceil\log_2(N+1)\rceil=O(\log N)$.

% Для варианта 2 в лучшем случае требуется $\lfloor\log_2N\rfloor$ шагов
% цикла и одна итоговая проверка, поэтому
% $T_{\min}(N)=\lfloor\log_2N\rfloor+1=O(\log N)$. В худшем случае
% $T_{\max}(N)=\lceil\log_2N\rceil+1=O(\log N)$.

% \section{Результаты сравнительного анализа}
% Результаты теоретической оценки ресурсной эффективности алгоритмов отображены в
% таблице~\ref{tab:comparison}.

% \begin{table}[H]
% 	\centering
% 	\caption{Сравнительный анализ ресурсной эффективности алгоритмов}
% 	\label{tab:comparison}
% 	\begin{tabularx}{\textwidth}{|X|c|c|X|>{\centering\arraybackslash}X|}
% 		\hline
% 		Реализация                           & $T_{\min}(N)$                                  & $T_{\max}(N)$ &
% 		Требования, накладываемые алгоритмом & Дополнительная память                                                         \\ \hline
% 		Линейный поиск
% 		                                     & $O(1)$                                         & $O(N)$        & Нет & $O(1)$ \\ \hline
% 		Бинарный поиск, нерекурсивный, вариант 1
% 		                                     & $O(1)$                                         & $O(\log N)$
% 		                                     & Массив должен быть предварительно отсортирован & $O(1)$                       \\ \hline
% 		Бинарный поиск, нерекурсивный, вариант 2
% 		                                     & $O(\log N)$                                    & $O(\log N)$
% 		                                     & Массив должен быть предварительно отсортирован & $O(1)$                       \\ \hline
% 		Бинарный поиск, рекурсивный
% 		                                     & $O(1)$                                         & $O(\log N)$
% 		                                     & Массив должен быть предварительно отсортирован & $O(\log N)$                  \\ \hline
% 	\end{tabularx}
% \end{table}

\section{Выводы}
% По результатам сравнительного анализа сделаны следующие выводы:
% \begin{itemize}
% 	\item бинарный поиск асимтотически эффективнее по времени чем линейный
% 	      поиск --- $O(\log_2N)$ против $O(N)$ в худшем случае, но требует
% 	      предварительной сортировки;
% 	\item бинарный поиск с ранним выходом и рекурсивный бинарный поиск
% 	      асимптотически эквивалентны по времени в лучшем и худшем случае, но не
% 	      рекурсивная реализация требует асимптотически меньше памяти;
% 	\item бинарный нерекурсивный поиск в двух вариантах различается только
% 	      асимптотикой лучшего случая --- $O(N)$ против $O(\log_2N)$;
% 	\item рекурсивная реализация асимптотически наихудшая по памяти среди
% 	      всех реализаций.
% \end{itemize}

\chapter*{ЗАКЛЮЧЕНИЕ}
% \addcontentsline{toc}{chapter}{ЗАКЛЮЧЕНИЕ}
% Цель работы достигнута: выполнен сравнительный анализ ресурсной эффективности
% алгоритмов. Первый вариант реализации нерекурсивного алгоритма имеет лучшую
% асимптотическую эффективность по времени в худшем и лучшем случае,
% а также лучшую асимптотическую эффективность по памяти. Рекурсивная реализация
% бинарного поиска имеет худшую асимптотическую эффективность по памяти.
% Реализация линейного поиска имеет худшую асимптотическую эффективность по времени.

% В ходе работы были решены поставленные задачи:
% \begin{itemize}
% 	\item описана решаемая задача согласно варианту;
% 	\item разработаны рекурсивный и нерекурсивные алгоритмы согласно варианту;
% 	\item описаны разработанные алгоритмы схемами алгоритмов;
% 	\item описаны требования к программному обеспечению;
% 	\item описаны средства реализации алгоритмов;
% 	\item реализованы разработанные алгоритмы;
% 	\item выполнено тестирование реализации алгоритмов;
% 	\item выполнена теоретическая оценка затрачиваемой памяти реализацией каждого алгоритма;
% 	\item оценена временная эффективность реализаций алгоритмов в худшем и лучшем
% 	      случае в терминах числа сравнений искомой величины с другими величинами.
% 	\item представлены результаты сравнительного анализа в формате таблицы.
% 	\item сделаны выводы из сравнительного анализа реализации разработанных
% 	      алгоритмов решения одной и той же задачи, заданной вариантом, по кри-
% 	      териям временной и ёмкостной эффективности, а также дополнительных
% 	      требований, налагаемых алгоритмом.
% \end{itemize}

\makebibliography

\appendix
\chapter*{Приложение А}
\addcontentsline{toc}{chapter}{Приложение А}

\lstinputlisting[
	language=C++,
	caption={Заголовочный файл matrix.hpp},
	label={lst:matrix-hpp}
]{../../source_code/lab_2/matrix.hpp}

\lstinputlisting[
	language=C++,
	caption={Файл matrix.cpp},
	label={lst:matrix-cpp}
]{../../source_code/lab_2/matrix.cpp}

\lstinputlisting[
	language=C++,
	caption={Заголовочный файл benchmark\_support.hpp},
	label={lst:benchmark-support-hpp}
]{../../source_code/lab_2/benchmark_support.hpp}

\lstinputlisting[
	language=C++,
	caption={Файл benchmark\_support.cpp},
	label={lst:benchmark-support-cpp}
]{../../source_code/lab_2/benchmark_support.cpp}

\lstinputlisting[
	language=C++,
	caption={Файл benchmark.cpp},
	label={lst:benchmark-cpp}
]{../../source_code/lab_2/benchmark.cpp}

\lstinputlisting[
	language=C++,
	caption={Файл main.cpp},
	label={lst:main-cpp}
]{../../source_code/lab_2/main.cpp}

\lstinputlisting[
	language=C++,
	caption={Файл test.cpp},
	label={lst:test-cpp}
]{../../source_code/lab_2/test.cpp}

\end{document}
